\documentclass[11pt,a4paper,twocolumn]{article}
\usepackage[utf8]{inputenc}
\usepackage[T1]{fontenc}
\usepackage{amsmath, amssymb, graphicx, booktabs, tikz}
\usepackage{microtype}
\usepackage[hyphens]{url}
\usetikzlibrary{arrows.meta, positioning, shapes.geometric}
\raggedbottom

\title{How humans solve indoor orienteering courses}
\author{}
\date{\today}

\begin{document}

\maketitle

\begin{abstract}
Algorithmic analysis of indoor orienteering courses models the venue as a
graph and finds routes by exhaustive search. A competitor reading a map
mid-race cannot explore the whole graph, holds only a handful of items in
working memory, and must commit in seconds --- yet solves within-floor
navigation almost instantly, at a glance. This article develops a cognitive
account of human route finding: within-floor regions are solved perceptually
as single units of reasoning, so the human's effective search space is the
condensed \emph{cluster graph} of floor regions connected by labeled stairs.
On that graph, humans plan bidirectionally from both controls, prefer near
and goal-directed stairs, and commit to the first connection found rather
than the best one. We analyze the challenges this strategy faces --- one-way
passages, discontinuous map segments, working-memory overflow --- and the
training methods that address them, and describe a deterministic simulation
that enables new difficulty measures and empirical tests.
\end{abstract}

\section{Introduction}

A companion article introduces a graph model of indoor orienteering venues:
nodes at decision-relevant positions, edges for corridors and stairs, one-way
constraints, and floor membership resolved from level polygons; the identity
of a route is its sequence of labeled stair crossings~\cite{framework}. On
this model, algorithms find the shortest route by Dijkstra search and can
enumerate \emph{every} distinct route between two controls, shortest first.

A human competitor holding the map does none of this. They cannot enumerate;
they often cannot even verify optimality of the route they choose. Still,
experienced indoor orienteers solve legs quickly and mostly well. Two
observations frame this article:

\begin{enumerate}
    \item \textbf{Within a floor, route finding is perceptual, not
          deliberative.} Given a start and a target on the same map segment,
          a quick glance is usually sufficient: the runner \emph{sees} a
          path through the corridors the way one sees the way out of a
          drawn maze, without item-by-item reasoning.
    \item \textbf{Across floors, route finding is deliberative and
          memory-bound.} Choosing which staircases to take, in which order,
          is slow, error-prone, and limited by how many intermediate steps
          the runner can hold in mind while also running.
\end{enumerate}

The venue-graph model explains why the boundary falls exactly there, what the
human's effective search problem is, and which strategies make it manageable.
Understanding this has three uses: it predicts where courses are hard (a
difficulty model should count \emph{mental} work, not graph edges); it
suggests concrete training methods; and it can be simulated, giving analysis
tools a ``human'' route to compare against the optimal one.

\section{Background}

Research on wayfinding in multi-level buildings shows that vertical
transitions dominate difficulty and that people adopt identifiable
strategies. H\"olscher et al.~\cite{holscher2006}, studying a complex
multi-level building, distinguish a \emph{central-point} strategy (route via
well-known parts), a \emph{direction} strategy (head toward the horizontal
position of the goal first), and a \emph{floor} strategy (get to the goal's
floor first); experienced participants favored the floor strategy.
Wiener et al.~\cite{wiener2009} classify wayfinding tasks by the knowledge
available; map-based indoor orienteering is unusual in giving complete
survey knowledge (the map) under severe time pressure.

Cognitive constraints are well established: working memory holds only about
four chunks~\cite{cowan2001} (the classical estimate was seven~\cite{miller1956}),
and expert advantage in domains like chess comes largely from \emph{chunking}
--- perceiving familiar configurations as single units~\cite{chase1973}.
Decision makers under time pressure \emph{satisfice}: they take the first
option that clears an acceptability threshold rather than optimizing
\cite{simon1956}. In orienteering specifically, expert cognition centers on
anticipatory map reading and strategies that keep attention free for
execution~\cite{eccles2002}.

Indoor orienteering maps add a further, unusual burden: each floor is printed
as a separate map segment, so a single staircase appears \emph{twice} --- once
on each floor's segment --- linked only by its label. Following a route across
floors therefore requires repeated visual search across segments, a task with
no outdoor counterpart.

\section{The human search problem on the venue graph}

\subsection{Floor regions are solved perceptually}

The algorithmic framework condenses every region reachable through ordinary
two-way corridors into a single \emph{cluster}, because corridors never
change route identity~\cite{framework}. The cognitive claim of this article
is that the same condensation happens perceptually in the human solver:

\begin{quote}
A floor region visible on one map segment is, for a competent map reader, a
solved subproblem --- a single chunk. Route finding within it is fast,
parallel, and nearly effortless; route finding \emph{between} regions is
slow, serial, and memory-bound.
\end{quote}

This is consistent with the maze-tracing character of within-floor planning
(a perceptual task at which humans are quick) and with chunking theory: the
corridor layout is processed as a spatial gestalt, not as discrete graph
nodes~\cite{chase1973}. It immediately explains the empirical shape of indoor
difficulty guidelines, which grade courses almost entirely by counts of floor
changes, stairs, and one-way passages --- never by corridor complexity alone
\cite{framework}.

The consequence is that \textbf{the human's effective search graph is the
cluster graph}: one node per floor region, one edge per labeled staircase or
one-way passage. A venue with thousands of map positions collapses, for the
planner, to a few dozen clusters and 10--20 stair edges. Human route finding
is search on this small graph under severe resource limits.

\subsection{Resource limits}

On the cluster graph the human solver is bounded in three ways:

\begin{itemize}
    \item \textbf{Working memory.} A planned route must be carried as a
          stair-label sequence (e.g.\ $dcccCA$: down D, down C three floors,
          up C, up A). At roughly four chunks~\cite{cowan2001}, sequences
          beyond 4--6 crossings cannot be held reliably while running unless
          they are re-chunked (repeated letters compress well: $ccc$ is one
          chunk) or re-read en route.
    \item \textbf{Exploration budget.} Each candidate stair considered costs
          a visual fixation on the map plus a cross-segment search for where
          the stair lands. Under race pressure a runner evaluates only a
          handful of candidates per decision, nothing like the systematic
          expansion of a best-first search.
    \item \textbf{No global comparison.} Comparing two complete multi-floor
          routes requires holding both --- typically beyond capacity. Hence
          route \emph{choice} degenerates to route \emph{finding}: the first
          plan that closes is taken~\cite{simon1956}.
\end{itemize}

\subsection{The observed strategy}
\label{sec:observed-strategy}

From these constraints, the strategy reported by practitioners and consistent
with the wayfinding literature follows almost by necessity:

\begin{enumerate}
    \item \textbf{Same-floor first.} Check whether the leg can be solved
          within the current cluster. If yes, run it immediately --- even
          when a stair route would be shorter. This is the floor strategy of
          H\"olscher et al.~\cite{holscher2006} taken to its extreme, and it
          is cheap because within-cluster solving is perceptual.
    \item \textbf{Bidirectional planning.} Otherwise, alternate between two
          easy local questions: \emph{which stairs can I reach from here?}
          (forward) and \emph{which stairs would drop me at the control?}
          (backward). Both are single-segment perceptual queries. Planning
          from both ends halves the depth that must be bridged --- the
          classic advantage of bidirectional search, discovered here not by
          algorithm design but by memory pressure.
    \item \textbf{Local heuristics for ordering.} Among candidate stairs,
          prefer the near one, and the one whose far end lies toward the
          goal (the direction strategy~\cite{holscher2006}).
    \item \textbf{First connection wins.} When the forward and backward
          fragments meet in a common floor region, the plan closes:
          concatenate the forward labels with the backward labels in reverse
          discovery order, and start running. No alternatives are examined.
\end{enumerate}

The result is satisficing --- settling for a plan that is good enough rather
than searching for the best one~\cite{simon1956} --- deterministic-ish, and
fast --- and systematically non-optimal in ways the algorithmic analysis can
measure (Section~\ref{sec:simulation}).

\subsection{Memory of previous routes}

The search just described does not start from scratch on every leg. Each
solved leg leaves behind reusable knowledge: which stair labels exist and
roughly where they sit, which regions a given staircase connects, and which
passages are one-way. A staircase already traversed is a known quantity ---
its far end no longer needs to be searched for on the other segment, so
evaluating it as a candidate costs almost nothing. In chunking terms
\cite{chase1973}, earlier routes convert parts of the cluster graph from raw
map content into pre-formed chunks, simultaneously widening the effective
exploration budget and cheapening each planned crossing. This predicts two
observable effects: later legs revisiting familiar stairs are planned faster
than structurally equivalent early legs, and runners should show a bias
toward stairs they have already used even when a slightly better unfamiliar
alternative exists --- a bias course setters can exploit in the same way as
the commitment traps of Section~\ref{sec:challenges}.

\section{Challenges}
\label{sec:challenges}

\subsection{One-way passages}
One-way passages break the symmetry that bidirectional planning relies on: a
passage usable in the forward fragment may be illegal when mentally walked
backward from the control, and a plan that ignores this discovers the error
only on the ground, where recovery is expensive (the passage cannot be
re-entered against its direction). The technique that addresses this is
\emph{directional discipline in the backward pass}: when asking ``which
stairs/passages would bring me here?'', only arrows pointing \emph{into} the
region count. This is exactly the point where novices' backward planning
fails, and course setters exploit it deliberately; the difficulty guidelines
gate one-way passages accordingly~\cite{framework}.

\subsection{Discontinuous map segments}
A staircase's two ends appear on different map segments, linked only by its
label. Every crossing added to a plan therefore costs a visual search: find
the label on the other segment, re-orient, and continue. Errors here ---
matching the wrong label, or the right label on the wrong segment --- corrupt
the whole plan. Practitioners mitigate this with \emph{label vocabulary}
(memorizing the venue's stair letters and their approximate positions during
the run-in or from earlier legs) and with map handling: thumb on the current
position, map folded to keep the two active segments adjacent.

The problem sharpens when the segments span multiple \emph{pages} --- for
large venues the map may be printed on both sides of the sheet, or on
several sheets. A route crossing a page boundary then involves segments that
cannot be viewed simultaneously at all: the runner must flip or swap the map
mid-plan, losing both the thumb anchor on the current position and any
visual context from the segment just left. Bidirectional planning suffers
particularly, since the forward and backward fragments may live on different
pages, forcing the working-memory-bound part of the plan to bridge a gap
that no fold can close. Crossings between pages are therefore markedly more
expensive than crossings between co-visible segments, and both course
setters and difficulty models should treat the page boundary as an
additional structural cost.

\subsection{Working-memory overflow}
Legs requiring many crossings exceed the 4-chunk budget. Copers use three
techniques: \textbf{compression} (encode the plan as the letter string
$dcccCA$ and rehearse it verbally --- repeated letters and familiar pairs
compress into single chunks~\cite{chase1973,cowan2001});
\textbf{plan-as-you-go} (commit only to the first crossing or two, and
re-plan at an intermediate region --- trading map-reading time while moving
for memory load); and \textbf{offloading to the map} (re-derive the remainder
at each floor change rather than storing it). Which mix is optimal depends on
running speed, map complexity, and the runner's rehearsal skill --- a genuine
individual-differences question the framework can test.

\subsection{Commitment traps}
Satisficing means the first closed plan wins, so a course setter can place an
attractive early connection that is far from optimal --- a \emph{bait}. The
countermeasure is a cheap \emph{sanity bound}: before running, compare the
plan's rough length or crossing count against the leg's straight-line
distance; a plan that feels disproportionate warrants one more look. Experts
appear to run this check habitually and almost instantly; it is a prime
candidate for explicit training (Section~\ref{sec:training}).

\subsection{Error recovery}
When execution diverges from the plan (a missed one-way, a wrong stair), the
runner must \emph{relocate}: re-establish position, then re-plan from there.
The same-floor-first strategy makes relocation cheaper than it looks --- the
current cluster is perceptually solvable once position is known --- so the
critical skill is recognizing \emph{that} one is off-plan early. Anticipatory
map reading, the hallmark of expert orienteers~\cite{eccles2002}, serves
exactly this: knowing what should appear next makes deviations easy to
notice.

\section{Training the technique}
\label{sec:training}

The analysis above decomposes indoor route finding into trainable components.
Deliberate practice --- isolated components, immediate feedback, graded
difficulty~\cite{ericsson1993} --- suggests the following program, ordered
from perceptual to strategic. The open toolchain~\cite{repo} supports several
of these directly, and a companion training application~\cite{training-app}
implements the program below as a set of deliberate-practice drills with
timing and immediate feedback.

\begin{enumerate}
    \item \textbf{Cluster reading.} Timed drills: given a map segment and two
          points, trace the within-floor route. Trains the perceptual solver
          that everything else builds on; speed and error rate are easy to
          measure.
    \item \textbf{Stair vocabulary.} Given a segment, name all stairs and
          their directions (up/down) as fast as possible; then, from memory,
          point to a named stair. Builds the label--position mapping that
          cross-segment searches depend on.
    \item \textbf{Label-sequence memory.} Rehearse routes as letter strings
          ($dcccCA$); recall them after a physical effort interval (memory
          under load differs sharply from memory at rest). Practice
          compression: find the chunking of a given string that survives a
          200\,m run.
    \item \textbf{Backward planning.} From the control, list every
          stair/passage that can legally reach it, one-way arrows respected.
          Do this on paper courses until it becomes automatic.
    \item \textbf{Bidirectional closure drills.} Full legs, planned under a
          countdown, answering only: forward reachable stairs, backward
          reaching stairs, first meeting region. The countdown forces
          satisficing; the drill trains making the \emph{first} plan a good
          one by improving candidate ordering.
    \item \textbf{Bait detection.} Review legs where the human-style plan and
          the optimal route diverge (the toolchain enumerates all routes and
          can compute both~\cite{repo}); learn the structural patterns ---
          attractive near stair, hidden one-way, shorter route via
          unintuitive direction --- and practice the sanity bound.
    \item \textbf{Post-race comparison.} After events, compare the route
          actually run (from split times or memory) against the enumerated
          route inventory: was the loss in planning (wrong route) or
          execution (right route, run badly)? This closes the deliberate-%
          practice feedback loop at the level of whole races.
\end{enumerate}

Course setters can invert the same list: each drill corresponds to a
difficulty knob (cluster complexity, stair count and naming, sequence depth,
one-way density, bait placement) that the graph model makes explicit and
countable.

\section{Simulating the human solver}
\label{sec:simulation}

The strategy of Section~3.3 is precise enough to implement. The analysis
toolchain~\cite{repo} contains a deterministic \emph{human route} algorithm
operating on the same venue graph as the optimal analyzers:

\begin{enumerate}
    \item \textbf{Same-floor short-circuit:} if the leg is solvable through
          unlabeled corridors alone (one-way corridors included in their
          legal direction), return that route immediately.
    \item \textbf{Alternating bidirectional search} on the cluster graph:
          one expansion per side per turn, forward side first. Candidate
          stairs from the current region are ordered nearest-first, with
          straight-line direction to the goal as tie-break; each side visits
          a region at most once (a human does not re-explore).
    \item \textbf{First valid meeting wins:} when a region is reached by
          both sides, forward labels are concatenated with the backward
          side's labels in reverse discovery order; a meeting whose two
          fragments would reuse the same physical staircase is rejected and
          the search continues. The concrete path is stitched through
          cached within-region shortest paths, and the search stops.
\end{enumerate}

The simulation is intentionally non-optimal, produces exactly one route, and
is deterministic, so its predictions are testable. It enables measures the
optimal analysis cannot provide: the \emph{human--optimal gap} (length ratio
between the simulated human route and the true shortest route --- a direct
bait detector); \emph{plan depth} (crossings in the plan, a working-memory
load estimate); and \emph{mental steps} (expansions before the plan closed, a
planning-time estimate). A leg can be easy for the simulation and hard for
Dijkstra-based scoring, or vice versa; systematic divergence between the two,
compared against split-time data, is a sharp empirical test of whether human
difficulty is better predicted by the human model. A companion plan in the
repository describes the intended refinements: exposing the exploration
trace for narration, parameterizing the cognitive limits (memory span,
exploration budget, heuristic weights), and fitting those parameters per
runner class against event data~\cite{repo}.

\section{Research directions}

\begin{enumerate}
    \item \textbf{Validate the same-floor bias:} from split times,
          measure how often runners take within-floor routes when a shorter
          stair route exists, by class and by time pressure.
    \item \textbf{Fit cognitive parameters:} treat memory span, exploration
          budget, and candidate-ordering weights as free parameters of the
          simulation; fit them to observed route choices per runner class.
    \item \textbf{Human-model difficulty scoring:} use plan depth, mental
          steps, and the human--optimal gap as terms in a difficulty model,
          and compare its predictive power against structural models on the
          same data~\cite{framework}.
    \item \textbf{Bait taxonomy:} classify legs with a large human--optimal
          gap by their structural cause (near-stair bait, one-way trap,
          direction illusion); this gives course setters a designed-%
          difficulty vocabulary.
    \item \textbf{Training studies:} test whether the drills of
          Section~\ref{sec:training} transfer --- e.g.\ does label-sequence
          memory training reduce time loss specifically on deep-sequence
          legs?
    \item \textbf{Instrumented stairwells:} place additional split-time
          units inside stairwells, not just at course controls, to measure
          time spent per floor change directly rather than inferring it, and
          to separate planning delay (deciding which stair to take) from
          execution delay (running it) within a single leg.
\end{enumerate}

\section{Conclusion}

Human indoor-orienteering route finding is bounded search on the condensed
cluster graph: floor regions are solved at a glance and treated as single
chunks, while stair-sequence planning is deliberative, bidirectional,
satisficing (Section~\ref{sec:observed-strategy}), and limited by a working
memory of a few chunks. The hard parts
of indoor navigation --- one-way passages, discontinuous map segments, deep
stair sequences, commitment traps --- are exactly the points where this
strategy's assumptions break, and each has an identifiable countermeasure and
a corresponding training drill. Because the strategy is precise enough to
simulate, human-style route finding can now be computed alongside optimal
routes, turning claims about ``what a runner would do'' into measurable,
falsifiable quantities.

\begin{thebibliography}{9}

\bibitem{framework}
\emph{A graph-based framework for quantitative analysis of navigational
complexity in indoor orienteering courses}. Companion article, same
repository.

\bibitem{holscher2006}
H\"olscher, C., Meilinger, T., Vrachliotis, G., Br\"osamle, M., \& Knauff, M.
(2006). Up the down staircase: Wayfinding strategies in multi-level
buildings. \emph{Journal of Environmental Psychology}, 26(4), 284--299.

\bibitem{wiener2009}
Wiener, J. M., B\"uchner, S. J., \& H\"olscher, C. (2009). Taxonomy of human
wayfinding tasks: A knowledge-based approach. \emph{Spatial Cognition \&
Computation}, 9(2), 152--165.

\bibitem{miller1956}
Miller, G. A. (1956). The magical number seven, plus or minus two: Some
limits on our capacity for processing information. \emph{Psychological
Review}, 63(2), 81--97.

\bibitem{cowan2001}
Cowan, N. (2001). The magical number 4 in short-term memory: A
reconsideration of mental storage capacity. \emph{Behavioral and Brain
Sciences}, 24(1), 87--114.

\bibitem{chase1973}
Chase, W. G., \& Simon, H. A. (1973). Perception in chess. \emph{Cognitive
Psychology}, 4(1), 55--81.

\bibitem{simon1956}
Simon, H. A. (1956). Rational choice and the structure of the environment.
\emph{Psychological Review}, 63(2), 129--138.

\bibitem{eccles2002}
Eccles, D. W., Walsh, S. E., \& Ingledew, D. K. (2002). A grounded theory of
expert cognition in orienteering. \emph{Journal of Sport and Exercise
Psychology}, 24(1), 68--88.

\bibitem{ericsson1993}
Ericsson, K. A., Krampe, R. T., \& Tesch-R\"omer, C. (1993). The role of
deliberate practice in the acquisition of expert performance.
\emph{Psychological Review}, 100(3), 363--406.

\bibitem{repo}
\emph{Indoor Orienteering Route Analysis} (source code, algorithm reference,
and datasets). \newblock
\url{github:objorke/indoor-o-route-analysis}.

\bibitem{training-app}
\emph{Indoor-O Training} (companion deliberate-practice training
application). \newblock
\url{github:objorke/indoor-o-training}.

\end{thebibliography}

\end{document}
